% !PRINTED 577
\weberpage{577}
%<<CONT
and from this one concludes again by Theorem 3. (applied to $\mathfrak p_1$)
\[ \Phi_1(t) \equiv P_1(t) \Phi_2(t) \pmod p. \]
In this the same consideration can be repeated which leads to the congruence
\[ \Phi_2(t) \equiv 0 \pmod{\mathfrak p_2 \ldots} \]
from which, again by 3.,
\[ \Phi_2(t) \equiv P_2(t) \Phi_3(t) \pmod p \]
is to be concluded. If one continues in this way until all the prime factors of $p$ have been taken into account, the following theorem results:

5. \gesperrt{If $\Phi(t)$ is an integral function in $R$ which becomes divisible by $p$ by the substitution $t=\tau$, then another integral function $\Phi_0(t)$ in $R$ can be determined so that}
\weq{29}{\Phi(t) \equiv \Phi_0(t) P(t) P_1(t) P_2(t) \ldots \pmod p}
\gesperrt{holds.}

\gesperrt{The product $P(t) P_1(t) P_2(t) \ldots$ is of degree $n = f+f_1+f_2+\cdots$ and is the integral rational form of lowest degree which becomes divisible by $p$ by the substitution $t=\tau$.}

Among the functions $\Phi(t)$ occurring in Theorem 5. there is also the integral function of degree $n$
\[ F(t) = N(t-\tau), \]
which vanishes for $t=\tau$. For this function $\Phi_0(t)$ becomes independent of $t$, and since both in $F(t)$ and in $P(t), P_1(t), P_2(t) \ldots$ the highest power of $t$ has the coefficient 1, it follows that $\Phi_0(t) = 1$. Hence we have the congruence
\weq{30}{N(t-\tau) \equiv P(t) P_1(t) P_2(t) \ldots \pmod p.}
If in the decomposition (25) of the prime $p$ we collect the equal prime factors into powers, then, if $\mathfrak p_1, \mathfrak p_2, \ldots$ are different prime ideals of the degrees $f_1, f_2, \ldots$, we can take the positive exponents $e_1, e_2, \ldots$ in such a way that
\weq{31}{p = \mathfrak p_1^{e_1} \mathfrak p_2^{e_2} \ldots,}
and
\weq{32}{n = e_1 f_1 + e_2 f_2 + \cdots}
hold. Formula (30) then assumes the shape
\weq{33}{N(t-\tau) \equiv [P_1(t)]^{e_1} [P_2(t)]^{e_2} \ldots \pmod p.}
